\chapter{Diagnóstico}
\label{chap:diagnostico}

Cuando aparece un error, conviene comprobar el sistema en el mismo orden seguido durante
la puesta en marcha: conexión PCIe, controlador, entorno de Python, modelo y datos.
La \cref{tab:troubleshooting} resume los fallos más habituales y dónde empezar a buscarlos.

\begin{longtable}{@{}L{0.25\textwidth}L{0.34\textwidth}L{0.33\textwidth}@{}}
\caption{Guía rápida de diagnóstico.}
\label{tab:troubleshooting}\\
\toprule
\textbf{Síntoma} & \textbf{Comprueba} & \textbf{Qué hacer} \\
\midrule
\endfirsthead

\toprule
\textbf{Síntoma} & \textbf{Comprueba} & \textbf{Qué hacer} \\
\midrule
\endhead

La tarjeta no aparece en \code{lspci}. &
Alimentación, cable FFC, montaje y configuración PCIe de la Raspberry Pi 5. &
Revisa primero la conexión física y la configuración PCIe. \\

La tarjeta aparece en PCIe, pero no existe \code{/dev/akida*}. &
Versión del kernel, headers, \code{dkms status}, \code{lsmod} y mensajes de \code{dmesg}. &
Revisa la instalación y compatibilidad del controlador. \\

Existe \code{/dev/akida*}, pero \code{akida.devices()} está vacío. &
Entorno virtual activo, versión de Akida y permisos del dispositivo. &
Comprueba que estás utilizando el entorno correcto y que el usuario puede acceder al dispositivo. \\

\code{import akida} produce un error. &
Versión de Python y paquetes instalados en el entorno activo. &
Activa el entorno virtual correcto o reinstala Akida dentro de él. \\

El archivo FBZ no carga. &
Ruta del archivo, versión de Akida y compatibilidad del modelo. &
Comprueba primero el modelo de prueba incluido con el repositorio. \\

El modelo no puede mapearse con \code{hw\_only=True}. &
Salida de \code{model.summary()}, capas utilizadas y versión del modelo. &
Comprueba que el modelo es compatible con la ejecución completa en la AKD1000. \\

La entrada es rechazada. &
Forma, tipo, rango y orden de los datos. &
Compara la entrada con \code{model.input\_shape} y con el preprocesado utilizado durante
la conversión. \\

La ejecución termina, pero la salida no es la esperada. &
Preprocesado, orden de las muestras y uso de \code{forward()} o \code{predict()}. &
Prueba primero una entrada cuyo resultado esperado sea conocido. \\

La latencia cambia mucho entre ejecuciones. &
Modo de reloj, modo de mapeo, calentamiento, carga del host y número de repeticiones. &
Repite la medida manteniendo fija la configuración. \\

No aparece una medida de potencia. &
Medición habilitada y estadísticas devueltas por Akida. &
Comprueba la configuración y aumenta la duración de la prueba si es necesario. \\

\bottomrule
\end{longtable}

\section{Información útil para diagnosticar un fallo}

La herramienta incluida con el repositorio puede recoger buena parte de la información
necesaria para localizar un problema:

\begin{lstlisting}[language=bash]
mkdir -p evidence
akd1000-bringup doctor --json > evidence/doctor.json
\end{lstlisting}

Si necesitas más detalle, guarda también el estado del sistema y del controlador:

\begin{lstlisting}[language=bash]
uname -a > evidence/uname.txt
lspci -nnk > evidence/lspci-nnk.txt
dkms status > evidence/dkms.txt
ls -l /dev/akida* > evidence/device-nodes.txt 2>&1
python -m pip freeze > evidence/pip-freeze.txt
sudo dmesg | grep -iE 'akida|edma|pcie' > evidence/driver-log.txt
\end{lstlisting}

Cuando informes de un problema, conserva también el comando que produjo el error y el
mensaje completo. Con esa información suele ser posible distinguir si el fallo está en
PCIe, el controlador, el entorno de Python, el modelo o los datos.